% ============================================================
%  CHAPTER 27: ATOMIC STRUCTURE AND SPECTRA
%  The Physics Codex: A Complete University Foundation
%  Korede Samuel Omotosho (Falcon)
%  Aurikrex Learning Systems, 2026
%
%  File: chapters/part5/ch27_atomic_structure.tex
% ============================================================

\chapter{Atomic Structure and Spectra}
\label{ch:atomic_structure}
\minitoc
\newpage

\chapterquote{If you thought that science was certain ---
well, that is just an error on your part.}{Richard Feynman}

\begin{objectivebox}
By the end of this chapter, you should be able to:
\begin{enumerate}[noitemsep]
  \item Describe the evidence for atomic structure from
  Rutherford's alpha-scattering experiment
  \item Explain the nuclear model of the atom and
  its limitations
  \item Describe Bohr's model of the hydrogen atom
  \item Calculate the energy levels and spectral
  line wavelengths of hydrogen
  \item Explain line emission and absorption spectra
  and their origins
  \item Define and apply the concept of the
  electron-volt
  \item Describe the photoelectric effect and
  explain why classical physics fails to account
  for it
  \item State and apply Einstein's photoelectric
  equation
  \item Describe wave-particle duality and de
  Broglie's hypothesis
  \item State the Heisenberg uncertainty principle
\end{enumerate}
\end{objectivebox}

% ============================================================
\section{Intuition: What is an Atom?}
% ============================================================

Ancient Greek philosophers proposed that matter could
be divided only so far before reaching an indivisible
unit --- the \textit{atomos}, meaning uncuttable. For
two thousand years this was philosophy, not physics.
Then, in the nineteenth century, chemists began to
compile evidence that matter behaved as though it
really was made of discrete, indivisible particles.
John Dalton formalised the atomic theory in 1808.

But the atom turned out to be far from simple. It has
structure. It contains smaller particles. It can be
understood only through quantum mechanics, a framework
of physics so radical that it forced physicists to
abandon their most basic intuitions about how nature
works.

This chapter traces the development of our
understanding of atomic structure: from Thomson's
plum-pudding model, to Rutherford's planetary model,
to Bohr's quantised orbits, to the modern quantum
mechanical picture. Each step was forced by
experimental evidence that the previous model could
not explain.

% ============================================================
\section{Models of the Atom}
% ============================================================

\subsection{Thomson's Plum-Pudding Model (1897)}

J.J. Thomson discovered the electron in 1897 by
studying cathode rays. He measured the charge-to-mass
ratio of electrons and showed they were far lighter
than any atom. Since atoms are electrically neutral
and contain electrons (negative), they must also
contain positive charge. Thomson proposed the
\textbf{plum-pudding model}: a sphere of diffuse
positive charge with electrons embedded in it like
plums in a pudding.

\subsection{Rutherford's Nuclear Model (1911)}

\begin{historybox}[Rutherford's Gold Foil Experiment, 1909--1911]
Ernest Rutherford (1871--1937) directed Hans Geiger
and Ernest Marsden to fire alpha particles at a thin
gold foil and observe where they went.

\textbf{Expected (plum-pudding model):} The positive
charge is spread diffusely throughout the atom. Alpha
particles, being much heavier than electrons, should
pass straight through with at most small deflections.

\textbf{Observed:}
\begin{itemize}[noitemsep]
  \item Most alpha particles passed straight through
  (as expected).
  \item A small fraction were deflected by large angles.
  \item A tiny fraction ($\sim 1$ in $8000$) bounced
  almost straight back.
\end{itemize}

Rutherford described his astonishment: ``It was as if
you fired artillery shells at tissue paper and they
came back and hit you.''

\textbf{Conclusion:} The positive charge and almost
all the mass must be concentrated in an extremely
small, dense \textbf{nucleus}. The electrons orbit
far outside. The atom is mostly empty space.
\end{historybox}

\begin{diagbox}[Rutherford's Gold Foil Experiment]
\begin{center}
\begin{tikzpicture}[scale=0.95, font=\small]

  % Detector screen is drawn first so trajectories and labels remain visible.
  \draw[thick, gray, dashed] (0,2.0) circle (4cm);

  % Alpha source
  \fill[gray!50] (-3.5,1.5) rectangle (-3.0,3.5);
  \draw[thick] (-3.5,1.5) rectangle (-3.0,3.5);
  \node[left, font=\small] at (-3.6,2.5)
    {$\alpha$ source};

  % Gold foil
  \fill[codexgold!40] (-0.1,-0.5) rectangle (0.1,5.5);
  \draw[very thick, codexgold] (-0.1,-0.5) -- (-0.1,5.5);
  \draw[very thick, codexgold] (0.1,-0.5) -- (0.1,5.5);
  \node[above, codexgold!80!black, fill=white, inner sep=1pt]
    at (0,5.6) {Gold foil};

  % Alpha beam (parallel rays from source; stop before the foil).
  \foreach \y in {1.8,2.0,2.2,2.5,3.2} {
    \draw[->, thick, codexred]
      (-3.0,\y) -- (-0.2,\y);
  }

  % Straight-through (most particles) reaches the detector screen.
  \draw[->, very thick, codexred!80]
    (0.2,2.5) -- (3.9686,2.5);
  \node[codexred!80, font=\small, fill=white, inner sep=1pt,
    align=center] at (2.5,1.95)
    {Most:\\straight through};

  % Small-angle deflection (away from the positive nucleus).
  \draw[->, very thick, codexred!80!black]
    (0.2,3.2) -- (3.4341,4.0511);

  % Large-angle deflections: close approaches bend away from the nucleus.
  \draw[->, very thick, codexorange!80!black]
    (0.2,2.2) -- (2.2133,5.3318);
  \draw[->, very thick, codexorange!80!black]
    (0.2,1.8) -- (2.8284,-0.8284);
  \node[above, font=\tiny, codexorange!80!black,
    fill=white, inner sep=1pt] at (2.1,5.6)
    {Large deflection};

  % Backscatter points back toward the source and reaches the screen.
  \draw[->, very thick, codexblue]
    (0.1,2.0) -- (-3.4395,4.0420);
  \node[above, codexblue, font=\small]
    at (-1.5,3.7) {Back-scatter\\($\sim$1 in 8000)};

  % Gold nucleus (tiny dot)
  \fill[codexgold] (0.0,2.0) circle (5pt);
  \draw[black!65, thin] (-0.7,0.65) -- (-0.05,1.85) -- (0.0,1.95);
  \node[black!75, font=\tiny] at (-1.0,0.5) {Nucleus};

  \node[black!75, font=\tiny] at (4.7,5.7)
    {Detector\\screen};

  % Conclusions box: shifted clear of the emerging rays and screen endpoints.
  \node[align=left, font=\small,
    draw=codexblue, rounded corners, fill=cyan!5,
    text width=5.8cm, inner sep=6pt] at (7.35,2.5)
    {\textbf{Conclusions:}\\
    $\bullet$ Most passes straight through\\
    \quad $\Rightarrow$ atom is mostly empty space\\
    $\bullet$ Few large deflections\\
    \quad $\Rightarrow$ nucleus is very small\\
    $\bullet$ Some back-scatter\\
    \quad $\Rightarrow$ nucleus is very dense\\
    \quad\quad and positively charged};

\end{tikzpicture}
\end{center}
\end{diagbox}

\begin{defbox}[The Nuclear Model of the Atom]
\begin{itemize}[noitemsep]
  \item The atom consists of a tiny, dense, positively
  charged \textbf{nucleus} at the centre.
  \item Electrons orbit the nucleus at relatively
  large distances (atomic radius $\sim 10^{-10}\ \text{m}$;
  nuclear radius $\sim 10^{-15}\ \text{m}$).
  \item The atom is therefore $\sim 10^5$ times larger
  than its nucleus in linear dimensions --- mostly
  empty space.
  \item The nucleus contains most of the atom's mass.
\end{itemize}

\textbf{Problem with the nuclear model:} Classical
electromagnetism predicts that a circular orbit
requires centripetal acceleration, and an accelerating
charge must radiate energy. The electron should spiral
into the nucleus in $\sim 10^{-11}\ \text{s}$. Atoms
are obviously stable. The model is incomplete.
\end{defbox}

% ============================================================
\section{Bohr's Model of Hydrogen}
% ============================================================

\begin{historybox}[Niels Bohr, 1913]
Niels Bohr (1885--1962) solved the stability problem
by a bold postulate: electrons can only exist in
certain specific orbits without radiating. These are
called \textbf{stationary states}. An electron in a
stationary state does not radiate despite its
acceleration. It can only emit or absorb radiation
when it \textit{jumps} from one orbit to another.
Bohr was essentially grafting quantum ideas onto
classical mechanics in a way he could not fully justify
--- but it worked astonishingly well for hydrogen.
\end{historybox}

\begin{lawbox}[Bohr's Postulates for the Hydrogen Atom]
\begin{enumerate}[noitemsep]
  \item Electrons orbit the nucleus in specific
  circular orbits called \textbf{stationary states}
  without radiating energy.
  \item The allowed orbits are those in which the
  angular momentum of the electron is an integer
  multiple of $h/(2\pi)$:
  \[
    mv_nr_n = n\frac{h}{2\pi} = n\hbar
    \quad (n = 1, 2, 3, \ldots)
  \]
  where $n$ is the \textbf{principal quantum number}.
  \item When an electron jumps from orbit $n_2$ to
  orbit $n_1$ (where $n_2 > n_1$), a photon is
  emitted with energy:
  \[
    hf = E_{n_2} - E_{n_1}
  \]
  When a photon of the right energy is absorbed, the
  electron jumps to a higher orbit.
\end{enumerate}
\end{lawbox}

\begin{processbox}[Energy Levels of Hydrogen]
Combining Bohr's quantisation condition with the
Coulomb force providing centripetal force gives the
energy of the $n$th orbit:
\[
  E_n = -\frac{13.6}{n^2}\ \text{eV}
  \quad (n = 1, 2, 3, \ldots)
\]
The ground state ($n=1$): $E_1 = -13.6\ \text{eV}$

First excited state ($n=2$): $E_2 = -3.40\ \text{eV}$

Second excited state ($n=3$): $E_3 = -1.51\ \text{eV}$

\textbf{Ionisation energy} = energy needed to remove
the electron from ground state = $+13.6\ \text{eV}$

The negative sign means the electron is bound to the
nucleus. Zero energy corresponds to an electron at
rest at infinity (just barely free).

Photon emitted when electron drops from $n_i$ to $n_f$:
\[
  hf = E_{n_i} - E_{n_f}
  = 13.6\left(\frac{1}{n_f^2} - \frac{1}{n_i^2}\right)
  \text{eV}
\]
\end{processbox}

\begin{definitionbox}[The Electron-Volt]
The \textbf{electron-volt} (eV) is the kinetic energy
gained by an electron when accelerated through a
potential difference of $1\ \text{V}$:
\[
  1\ \text{eV} = 1.6\times10^{-19}\ \text{J}
\]
Useful for atomic-scale energies. Convert to joules
when using SI equations ($E = hf$, $E = kT$, etc.).
\end{definitionbox}

\begin{diagbox}[Energy Level Diagram of Hydrogen and Spectral Series]
\begin{center}
\begin{tikzpicture}[scale=1.0, font=\small]

  % Energy levels are schematic; vertical spacings are not to scale.
  % n=1: -13.6 eV; n=2: -3.4; n=3: -1.51; n=4: -0.85; n=5: -0.54; n=\infty: 0

  % Draw levels
  \foreach \n/\E/\y in {
    1/{-13.6}/{0},
    2/{-3.40}/{3.5},
    3/{-1.51}/{4.6},
    4/{-0.85}/{5.2},
    5/{-0.54}/{5.55}} {
    \draw[thick] (1.5,\y) -- (5.5,\y);
    \node[left, font=\small] at (1.4,\y)
      {$n=\n$};
    \node[right, font=\small] at (9.0,\y)
      {$E=\E$ eV};
  }

  % The ionisation limit is the n = infinity level.
  \draw[thick] (1.5,6.2) -- (5.5,6.2);
  \node[left, font=\small] at (1.4,6.2) {$n=\infty$};
  \node[right, font=\small] at (9.0,6.2) {$E=0$ eV};

  % Ionisation label
  \node[above, codexblue, font=\small]
    at (3.5,6.3) {$E=0$: Ionisation limit};

  % === LYMAN SERIES (UV): drops to n=1 ===
  % n=2 to n=1
  \draw[->, very thick, codexblue]
    (2.0,3.5) -- (2.0,0.1);
  % n=3 to n=1
  \draw[->, thick, codexblue!80]
    (2.3,4.6) -- (2.3,0.1);
  % n=4 to n=1
  \draw[->, thick, codexblue!60]
    (2.6,5.2) -- (2.6,0.1);
  \node[left, codexblue, font=\small, align=right]
    at (1.3,2.0) {Lyman\\series\\(UV)};

  % === BALMER SERIES (visible/UV): drops to n=2 ===
  % n=3 to n=2
  \draw[->, very thick, codexred]
    (3.5,4.6) -- (3.5,3.6);
  % n=4 to n=2
  \draw[->, thick, codexred!80]
    (3.8,5.2) -- (3.8,3.6);
  % n=5 to n=2
  \draw[->, thick, codexred!60]
    (4.1,5.55) -- (4.1,3.6);
  \node[right, codexred, font=\small, align=left]
    at (6.0,4.0) {Balmer\\series\\(visible/UV)};

  % === PASCHEN SERIES (IR): drops to n=3 ===
  \draw[->, very thick, codexgreen]
    (4.7,5.2) -- (4.7,4.7);
  \draw[->, thick, codexgreen!80]
    (5.0,5.55) -- (5.0,4.7);
  \node[right, codexgreen, font=\small]
    at (6.0,5.4) {Paschen\\(IR)};

  % Visible light annotation for Balmer
  \draw[codexred, thick, decorate,
    decoration={brace, mirror, amplitude=4pt}]
    (3.4,3.5) -- (3.4,4.6);
  \node[right, codexred, font=\tiny, align=left]
    at (3.85,3.1) {H$\alpha$\\656 nm};

  % Make the qualitative vertical layout explicit for readers.
  \node[font=\scriptsize, text=black!65, align=center]
    at (7.3,2.2) {Schematic; vertical spacings\\not to scale};

\end{tikzpicture}
\end{center}
\small Each series corresponds to transitions ending
at a particular level. The Balmer series falls partly
in the visible spectrum (400--700 nm) and is the most
commonly observed in astronomy.
\end{diagbox}

\begin{examplebox}[27.1]
\stepcounter{workedexample}
Calculate the wavelength of the photon emitted when
a hydrogen electron transitions from $n = 3$ to
$n = 2$ (the H$\alpha$ line, the first line of the
Balmer series).\\
($h = 6.63\times10^{-34}\ \text{J\,s}$,
$c = 3.0\times10^8\ \text{m\,s}^{-1}$,
$1\ \text{eV} = 1.6\times10^{-19}\ \text{J}$)
\end{examplebox}

\begin{solutionbox}
Energy difference:
\[
  \Delta E = E_3 - E_2 = -1.51 - (-3.40)
  = 1.89\ \text{eV}
  = 1.89 \times 1.6\times10^{-19}
  = 3.024\times10^{-19}\ \text{J}
\]

Using $E = hf = hc/\lambda$:
\[
  \lambda = \frac{hc}{\Delta E}
  = \frac{6.63\times10^{-34} \times 3.0\times10^8}
    {3.024\times10^{-19}}
  = \frac{1.989\times10^{-25}}{3.024\times10^{-19}}
  = 6.58\times10^{-7}\ \text{m}
  = 658\ \text{nm}
\]

This is in the red part of the visible spectrum.
The accepted value is $656\ \text{nm}$ ---
excellent agreement. The H$\alpha$ line is the
strong red emission seen in photographs of hydrogen
nebulae across the galaxy.
\end{solutionbox}

% ============================================================
\section{Emission and Absorption Spectra}
% ============================================================

\begin{defbox}[Line Emission and Absorption Spectra]
\textbf{Line emission spectrum:} When a gas is heated
or excited electrically, atoms are excited to higher
energy levels. As electrons fall back to lower levels,
photons of specific wavelengths are emitted. When the
emitted light is passed through a prism or diffraction
grating, a series of \textbf{bright coloured lines}
on a dark background is observed. Each element has a
unique set of spectral lines --- a fingerprint.

\textbf{Line absorption spectrum:} When white light
passes through a cool gas, the atoms absorb photons
of the same specific wavelengths they would emit.
The transmitted light shows \textbf{dark lines}
(Fraunhofer lines) at the same wavelengths as the
emission spectrum.

\textbf{Continuous spectrum:} A hot solid or dense
gas emits light of all wavelengths (a rainbow).
This is the blackbody spectrum from Chapter 16.
\end{defbox}

\begin{diagbox}[Emission and Absorption Spectra Compared]
\begin{center}
\begin{tikzpicture}[scale=0.95, font=\small]

  % Smooth spectrum built from adjoining color ramps.

  % === CONTINUOUS SPECTRUM ===
  \begin{scope}[shift={(0,5)}]
    \node[font=\bfseries] at (4.0,1.6)
      {Continuous spectrum (hot solid)};
    % Smooth continuum rather than discrete-looking color blocks.
    \shade[left color=violet!60, right color=blue!60] (0,0) rectangle (0.775,1.0);
    \shade[left color=blue!60, right color=cyan!60] (0.775,0) rectangle (1.55,1.0);
    \shade[left color=cyan!60, right color=green!60] (1.55,0) rectangle (2.325,1.0);
    \shade[left color=green!60, right color=codexgold!60] (2.325,0) rectangle (3.1,1.0);
    \shade[left color=codexgold!60, right color=orange!60] (3.1,0) rectangle (3.875,1.0);
    \shade[left color=orange!60, right color=codexred!60] (3.875,0) rectangle (4.65,1.0);
    \shade[left color=codexred!60, right color=codexred!60] (4.65,0) rectangle (5.425,1.0);
    \shade[left color=codexred!60, right color=codexred!60] (5.425,0) rectangle (6.2,1.0);
    \draw[thick] (0,0) rectangle (6.2,1.0);
    \node[below, font=\tiny] at (0.5,-0.1) {400 nm};
    \node[below, font=\tiny] at (5.7,-0.1) {700 nm};
    \node[below, font=\tiny] at (3.1,-0.1) {Wavelength (nm)};
  \end{scope}

  % === EMISSION SPECTRUM ===
  \begin{scope}[shift={(0,2.5)}]
    \node[font=\bfseries] at (4.0,1.6)
      {Emission spectrum (hot hydrogen gas)};
    % Dark background
    \fill[black] (0,0) rectangle (6.2,1.0);
    \draw[thick] (0,0) rectangle (6.2,1.0);
    % Balmer line positions use x=(lambda-400)/300*6.2.
    % H-alpha 656 nm (red)
    \draw[very thick, codexred]
      (5.29,0) -- (5.29,1.0);
    % H-beta 486 nm (blue-green)
    \draw[very thick, cyan!80]
      (1.78,0) -- (1.78,1.0);
    % H-gamma 434 nm (violet)
    \draw[very thick, blue!70]
      (0.70,0) -- (0.70,1.0);
    % H-delta 410 nm (violet)
    \draw[very thick, violet]
      (0.21,0) -- (0.21,1.0);
    % Labels
    \node[above, codexred, font=\tiny] at (5.29,1.05)
      {656};
    \node[above, cyan, font=\tiny] at (1.78,1.05)
      {486};
    \node[above right, blue, font=\tiny] at (0.70,1.05)
      {434};
    \node[above left, violet, font=\tiny] at (0.21,1.05)
      {410};
  \end{scope}

  % === ABSORPTION SPECTRUM ===
  \begin{scope}[shift={(0,0)}]
    \node[font=\bfseries] at (4.0,1.6)
      {Absorption spectrum (cool hydrogen)};
    % Continuous rainbow background for absorption.
    \shade[left color=violet!60, right color=blue!60] (0,0) rectangle (0.775,1.0);
    \shade[left color=blue!60, right color=cyan!60] (0.775,0) rectangle (1.55,1.0);
    \shade[left color=cyan!60, right color=green!60] (1.55,0) rectangle (2.325,1.0);
    \shade[left color=green!60, right color=codexgold!60] (2.325,0) rectangle (3.1,1.0);
    \shade[left color=codexgold!60, right color=orange!60] (3.1,0) rectangle (3.875,1.0);
    \shade[left color=orange!60, right color=codexred!60] (3.875,0) rectangle (4.65,1.0);
    \shade[left color=codexred!60, right color=codexred!60] (4.65,0) rectangle (5.425,1.0);
    \shade[left color=codexred!60, right color=codexred!60] (5.425,0) rectangle (6.2,1.0);
    \draw[thick] (0,0) rectangle (6.2,1.0);
    % Narrow dark lines centered at the same positions as emission
    \fill[black] (5.26,0) rectangle (5.32,1.0);
    \fill[black] (1.75,0) rectangle (1.81,1.0);
    \fill[black] (0.67,0) rectangle (0.73,1.0);
    \fill[black] (0.18,0) rectangle (0.24,1.0);
    % Labels
    \node[below, font=\tiny] at (5.29,-0.1) {656};
    \node[below, font=\tiny] at (1.78,-0.1) {486};
    \node[below right, font=\tiny] at (0.70,-0.1) {434};
    \node[below left, font=\tiny] at (0.21,-0.1) {410};
  \end{scope}

  % Key observation
  \node[align=center, font=\small, codexblue]
    at (4.0,-1.0)
    {Same wavelengths appear as bright lines
    in emission and dark lines in absorption.\\
    This is used in stellar spectroscopy to identify
    elements in stars.};

\end{tikzpicture}
\end{center}
\end{diagbox}

% ============================================================
\section{The Photoelectric Effect}
% ============================================================

\begin{processbox}[The Photoelectric Effect: Observations]
When light of sufficiently high frequency shines on a
metal surface, electrons are emitted from the surface.
This is the \textbf{photoelectric effect}.

Key experimental observations:
\begin{enumerate}[noitemsep]
  \item Electrons are only emitted if the frequency of
  the light exceeds a minimum value called the
  \textbf{threshold frequency} $f_0$. Below $f_0$,
  no electrons are emitted, however bright the light.
  \item Above $f_0$, electrons are emitted
  \textit{immediately}, even in very faint light. There
  is no time delay.
  \item The maximum kinetic energy of the emitted
  electrons depends only on the \textit{frequency}
  of the light, not on its intensity.
  \item Increasing the intensity (brightness) of the
  light above $f_0$ increases the \textit{number} of
  emitted electrons per second, but not their energy.
\end{enumerate}

\textbf{Why classical wave physics fails:}
Classical wave theory predicts that (1) any frequency
should work if the light is bright enough, (2) electrons
should take time to accumulate enough energy, and (3)
brighter light should produce more energetic electrons.
All three predictions are wrong. Classical theory
completely fails to explain the photoelectric effect.
\end{processbox}

\begin{historybox}[Einstein's Explanation, 1905]
Albert Einstein (1879--1955) explained the photoelectric
effect by proposing that light comes in discrete packets
of energy called \textbf{photons}. Each photon carries
energy $E = hf$ (where $h$ is Planck's constant). This
was Einstein's Nobel Prize paper (1921), not relativity.

The explanation: One photon interacts with one electron.
If the photon energy $hf$ is less than the work
function $\phi$ (the minimum energy needed to escape the
metal surface), no electron can escape. If $hf > \phi$,
the electron escapes with kinetic energy
$KE_{max} = hf - \phi$. Brighter light means more
photons per second, so more electrons per second --- but
each photon still has the same energy, so the maximum
KE is unchanged.
\end{historybox}

\begin{lawbox}[Einstein's Photoelectric Equation]
\[
  hf = \phi + \frac{1}{2}mv_{max}^2
\]
or equivalently:
\[
  KE_{max} = hf - \phi
\]
where:
\begin{itemize}[noitemsep]
  \item $h = 6.63\times10^{-34}\ \text{J\,s}$
  (Planck's constant)
  \item $f$ = frequency of incident light (Hz)
  \item $\phi = hf_0$ = work function of the metal (J or eV)
  \item $f_0$ = threshold frequency
  \item $KE_{max} = \frac{1}{2}mv_{max}^2$ = maximum
  kinetic energy of emitted electrons
\end{itemize}

\textbf{Stopping potential} $V_s$: the minimum
reverse potential needed to stop all emitted electrons:
\[
  eV_s = KE_{max} = hf - \phi
\]
\end{lawbox}

\begin{diagbox}[Photoelectric Effect: Experimental Setup and Graph]
\begin{center}
\begin{tikzpicture}[scale=0.9, font=\small]

  % === APPARATUS ===
  \begin{scope}[shift={(0,0)}]
    % Vacuum tube
    \fill[cyan!5] (0.3,0.3) rectangle (5.7,4.7);
    \draw[thick] (0.3,0.3) rectangle (5.7,4.7);
    \node[black!80, font=\small] at (3.0,4.2)
      {Evacuated tube};

    % Photocathode (metal surface)
    \fill[gray!50] (0.3,1.0) rectangle (1.0,4.0);
    \draw[thick] (0.3,1.0) rectangle (1.0,4.0);
    \node[font=\tiny, align=right, anchor=north east] at (0.92,0.05)
      {Metal\\cathode};

    % Anode
    \fill[gray!30] (4.8,1.5) rectangle (5.3,3.5);
    \draw[thick] (4.8,1.5) rectangle (5.3,3.5);
    \node[font=\fontsize{4.5}{5}\selectfont, anchor=east] at (4.9,1.3)
      {Anode};

    % Incident light
    \foreach \y in {1.5,2.0,2.5,3.0,3.5} {
      \draw[->, thick, codexgold]
        (-0.8,\y) -- (0.2,\y);
    }
    \node[left, codexgold] at (-0.9,2.5)
      {Light\\($f > f_0$)};

    % Emitted electrons
    \draw[->, thick, codexred]
      (1.1,2.0) -- (4.7,2.5);
    \draw[->, thick, codexred]
      (1.1,2.5) -- (4.7,2.5);
    \draw[->, thick, codexred]
      (1.1,3.0) -- (4.7,2.5);
    \node[above, codexred, font=\tiny]
      at (2.9,2.8) {Electrons};

    % Series ammeter and adjustable retarding supply; both electrodes are wired.
    \draw[thick] (1.0,1.0) -- (1.0,-0.5) -- (3.2,-0.5);
    \fill[white] (3.2,-0.8) rectangle (3.8,-0.2);
    \draw[thick] (3.2,-0.8) rectangle (3.8,-0.2);
    \node[font=\tiny] at (3.5,-0.5) {G};
    \draw[thick] (3.8,-0.5) -- (4.0,-0.5);
    \draw[very thick] (4.0,-0.82) -- (4.0,-0.18);
    \draw[very thick] (4.3,-0.70) -- (4.3,-0.30);
    \node[above, font=\tiny] at (4.0,-0.16) {$+$};
    \node[above, font=\tiny] at (4.3,-0.16) {$-$};
    \draw[thick] (4.3,-0.5) -- (5.0,-0.5) -- (5.0,1.5);
    \node[below, font=\tiny, align=center] at (4.55,-0.75)
      {Adjustable reverse bias $V_s$};
  \end{scope}

  % === KE vs FREQUENCY GRAPH ===
  \begin{axis}[
    at={(200pt,0pt)},
    width=7cm, height=6cm,
    xlabel={Frequency $f$ ($10^{14}\,\mathrm{Hz}$)},
    ylabel={$KE_{max}$ (eV)},
    xmin=0, xmax=10,
    ymin=-1.5, ymax=3.2,
    grid=major, grid style={gray!15},
    xtick={0,2,3,4,6,8,10},
    xticklabels={0,2,3,4,6,8,10},
    ytick={-1,0,1,2,3},
    yticklabels={$-1$,$0$,$1$,$2$,$3$},
  ]
    % f0=3e14 Hz, phi=1.24 eV; slope h=4.136e-15 eV s.
    \addplot[codexblue, very thick, domain=3:10]
      {0.4136*(x-3)};
    % Below threshold: KE = 0 (no emission)
    \addplot[gray, dashed, domain=0:3]{0};
    % Mark threshold
    \fill[codexred] (axis cs:3,0) circle (3pt);
    \node[above right, codexred, font=\tiny, align=left,
      text width=1.9cm]
      at (axis cs:3.1,1.25) {$f_0=3\times10^{14}\,\mathrm{Hz}$ threshold};
    % Work function annotation
    \draw[<->, codexgold!55!black, thick]
      (axis cs:0.12,-1.43) -- (axis cs:3,-1.43);
    \node[above, codexgold!55!black, font=\tiny, align=center, text width=1.4cm]
      at (axis cs:2.8,-1.25) {$\phi=hf_0\approx1.24$ eV};
    % Extension of line to y-intercept
    \addplot[codexblue, dashed, domain=0:3]
      {0.4136*(x-3)};
    \fill[codexgold!55!black] (axis cs:0,-1.2408) circle (3pt);
    \node[above right, codexgold!55!black, font=\tiny, align=left,
      text width=2.2cm, fill=white, inner sep=1pt]
      at (axis cs:2.6,-0.65) {$y$-intercept $=-\phi\approx-1.24$ eV};
    \node[black!80, font=\tiny] at (axis cs:1.4,0.25) {No emission};
  \end{axis}

\end{tikzpicture}
\end{center}
\small The line follows $KE_{\max}=hf-\phi$. With the
frequency axis scaled in $10^{14}\,\mathrm{Hz}$, its plotted
slope is $0.414\,\mathrm{eV}/(10^{14}\,\mathrm{Hz})$,
equivalent to $h=4.136\times10^{-15}\,\mathrm{eV\,s}$.
\end{diagbox}

\begin{examplebox}[27.2]
\stepcounter{workedexample}
Ultraviolet light of wavelength $250\ \text{nm}$
falls on a zinc surface with work function
$4.3\ \text{eV}$.\\
(a) Calculate the energy of each photon in eV.\\
(b) Calculate the maximum kinetic energy of the
emitted electrons.\\
(c) Calculate the maximum speed of the emitted
electrons.\\
(d) Calculate the stopping potential.\\
($h = 6.63\times10^{-34}\ \text{J\,s}$,
$c = 3.0\times10^8\ \text{m\,s}^{-1}$,
$m_e = 9.11\times10^{-31}\ \text{kg}$,
$e = 1.6\times10^{-19}\ \text{C}$)
\end{examplebox}

\begin{solutionbox}
\textbf{(a)} Photon energy:
\[
  E = \frac{hc}{\lambda}
  = \frac{6.63\times10^{-34} \times 3.0\times10^8}
    {250\times10^{-9}}
  = \frac{1.989\times10^{-25}}{2.5\times10^{-7}}
  = 7.956\times10^{-19}\ \text{J}
  = \frac{7.956\times10^{-19}}{1.6\times10^{-19}}
  = 4.97\ \text{eV}
\]

\textbf{(b)} Maximum kinetic energy:
\[
  KE_{max} = hf - \phi = 4.97 - 4.3
  = 0.67\ \text{eV}
  = 0.67 \times 1.6\times10^{-19}
  = 1.072\times10^{-19}\ \text{J}
\]

\textbf{(c)} Maximum speed:
\[
  \frac{1}{2}m_e v_{max}^2 = 1.072\times10^{-19}
\]
\[
  v_{max} = \sqrt{\frac{2 \times 1.072\times10^{-19}}
    {9.11\times10^{-31}}}
  = \sqrt{2.354\times10^{11}}
  = 4.85\times10^5\ \text{m\,s}^{-1}
\]

\textbf{(d)} Stopping potential:
\[
  V_s = \frac{KE_{max}}{e}
  = \frac{1.072\times10^{-19}}{1.6\times10^{-19}}
  = 0.67\ \text{V}
\]
\end{solutionbox}

% ============================================================
\section{Wave-Particle Duality}
% ============================================================

The photoelectric effect demonstrates that light,
which we know is a wave (it diffracts, interferes,
etc.), also behaves as discrete particles (photons).
This is wave-particle duality for light.

In 1924, Louis de Broglie proposed that matter
(which we know consists of particles) also has
wave properties.

\begin{lawbox}[De Broglie's Hypothesis]
Every particle with momentum $p$ has an associated
\textbf{de Broglie wavelength}:
\[
  \lambda = \frac{h}{p} = \frac{h}{mv}
\]
where $h$ is Planck's constant, $m$ is the particle's
mass and $v$ its speed.

This was confirmed experimentally by Davisson and
Germer in 1927: a beam of electrons was diffracted
by a crystal lattice, producing an interference
pattern exactly like X-rays of the same wavelength.
Particles diffract --- they have wave properties.
\end{lawbox}

\begin{examplebox}[27.3]
\stepcounter{workedexample}
Calculate the de Broglie wavelength of:
(a) An electron accelerated through $1000\ \text{V}$.\\
(b) A cricket ball of mass $0.16\ \text{kg}$ moving
at $30\ \text{m\,s}^{-1}$.\\
Comment on whether wave behaviour is observable
in each case.\\
($h = 6.63\times10^{-34}\ \text{J\,s}$,
$m_e = 9.11\times10^{-31}\ \text{kg}$,
$e = 1.6\times10^{-19}\ \text{C}$)
\end{examplebox}

\begin{solutionbox}
\textbf{(a)} Electron accelerated through $1000\ \text{V}$:

Kinetic energy gained: $KE = eV = 1.6\times10^{-19}
\times 1000 = 1.6\times10^{-16}\ \text{J}$

Speed: $v = \sqrt{2KE/m_e} = \sqrt{2 \times 1.6\times10^{-16}
/ 9.11\times10^{-31}} = 1.876\times10^7\ \text{m\,s}^{-1}$

Momentum: $p = m_e v = 9.11\times10^{-31}
\times 1.876\times10^7 = 1.709\times10^{-23}\ \text{kg\,m\,s}^{-1}$

De Broglie wavelength:
\[
  \lambda = \frac{h}{p}
  = \frac{6.63\times10^{-34}}{1.709\times10^{-23}}
  = 3.88\times10^{-11}\ \text{m} = 0.0388\ \text{nm}
\]
This is comparable to atomic spacings ($\sim 0.1\ \text{nm}$),
so electron diffraction from crystals is readily
observed. Electron microscopes exploit this.

\textbf{(b)} Cricket ball:
\[
  p = mv = 0.16 \times 30 = 4.8\ \text{kg\,m\,s}^{-1}
\]
\[
  \lambda = \frac{6.63\times10^{-34}}{4.8}
  = 1.38\times10^{-34}\ \text{m}
\]
This is $10^{19}$ times smaller than a proton
($\sim 10^{-15}\ \text{m}$). Completely undetectable.
Macroscopic objects have de Broglie wavelengths so
small that quantum wave effects are entirely
negligible --- this is why classical physics works
for everyday objects.
\end{solutionbox}

% ============================================================
\section{The Heisenberg Uncertainty Principle}
% ============================================================

\begin{lawbox}[Heisenberg Uncertainty Principle]
It is impossible to simultaneously know both the
exact position and the exact momentum of a particle.
The uncertainties are related by:
\[
  \Delta x \cdot \Delta p \geq \frac{h}{4\pi}
  \approx \frac{\hbar}{2}
\]
where $\Delta x$ is the uncertainty in position and
$\Delta p$ is the uncertainty in momentum.

Similarly, for energy and time:
\[
  \Delta E \cdot \Delta t \geq \frac{h}{4\pi}
\]

This is not a limitation of our instruments ---
it is a fundamental feature of nature. A particle
does not \textit{have} a precise position and precise
momentum simultaneously. The act of measuring one
disturbs the other.

This is why the Bohr model, which assumes precise
electron orbits, is ultimately incomplete. Electrons
do not orbit in fixed paths --- they exist as
probability clouds (orbitals) described by quantum
mechanics.
\end{lawbox}

% ============================================================
\section{Chapter Summary}
% ============================================================

\begin{summarybox}
\textbf{Thomson:} plum-pudding model (electrons
embedded in diffuse positive charge).

\textbf{Rutherford's scattering experiment:}
most $\alpha$ pass through; some deflect greatly;
rare back-scatter $\Rightarrow$ tiny, dense, positive nucleus;
atom mostly empty space.

\textbf{Bohr model:} electrons in fixed orbits;
quantised angular momentum; energy levels of H:
$E_n = -13.6/n^2$ eV.
Photon emitted/absorbed on transition:
$hf = E_{n_2} - E_{n_1}$.

\textbf{1 eV $= 1.6\times10^{-19}$ J}

\textbf{Spectral lines:} each element has unique
emission/absorption spectrum (fingerprint).
Balmer series (visible): transitions to $n=2$.

\textbf{Photoelectric effect:}
$hf = \phi + KE_{max}$;
threshold frequency $f_0 = \phi/h$.
Proves light is quantised (photons).

\textbf{Photon energy:} $E = hf = hc/\lambda$

\textbf{De Broglie wavelength:} $\lambda = h/p = h/(mv)$
Particles have wave properties; waves have particle
properties --- wave-particle duality.

\textbf{Heisenberg:}
$\Delta x\,\Delta p \geq h/4\pi$
(position and momentum cannot both be precisely known).
\end{summarybox}

% ============================================================
\newpage
\begin{exercisebox}[27]

\subsection*{Section A: Multiple Choice}

\textbf{A1.} Rutherford's alpha-scattering experiment
showed that:

\mcoption{A}{Atoms contain electrons}
\mcoption{B}{Most of the atom's mass is spread
uniformly through its volume}
\mcoption{C}{The atom has a tiny, dense, positively
charged nucleus surrounded by mostly empty space}
\mcoption{D}{The atom is made of protons only}
\ansref{Ch27-A1}

\vspace{0.4cm}

\textbf{A2.} According to Bohr's model, the energy
of the $n=4$ level of hydrogen is:

\mcoption{A}{$-3.40\ \text{eV}$}
\mcoption{B}{$-1.51\ \text{eV}$}
\mcoption{C}{$-0.85\ \text{eV}$}
\mcoption{D}{$-0.54\ \text{eV}$}
\ansref{Ch27-A2}

\vspace{0.4cm}

\textbf{A3.} The minimum frequency of light needed
to cause photoelectric emission from a metal surface is:

\mcoption{A}{The resonant frequency}
\mcoption{B}{The threshold frequency}
\mcoption{C}{The Lyman frequency}
\mcoption{D}{The Balmer frequency}
\ansref{Ch27-A3}

\vspace{0.4cm}

\textbf{A4.} In the photoelectric effect, increasing
the intensity of light above the threshold frequency
increases:

\mcoption{A}{The maximum kinetic energy of
photoelectrons}
\mcoption{B}{The threshold frequency}
\mcoption{C}{The number of photoelectrons emitted
per second}
\mcoption{D}{The work function of the metal}
\ansref{Ch27-A4}

\vspace{0.4cm}

\textbf{A5.} The de Broglie wavelength of a particle
is:

\mcoption{A}{$\lambda = hf$}
\mcoption{B}{$\lambda = h/p$}
\mcoption{C}{$\lambda = p/h$}
\mcoption{D}{$\lambda = hv$}
\ansref{Ch27-A5}

\vspace{0.4cm}

\textbf{A6.} An electron transition from $n=4$
to $n=2$ in hydrogen produces a photon of energy:

\mcoption{A}{$0.66\ \text{eV}$}
\mcoption{B}{$2.55\ \text{eV}$}
\mcoption{C}{$3.40\ \text{eV}$}
\mcoption{D}{$4.91\ \text{eV}$}
\ansref{Ch27-A6}

\vspace{0.4cm}

\textbf{A7.} The stopping potential in a
photoelectric experiment measures:

\mcoption{A}{The work function of the metal}
\mcoption{B}{The threshold frequency}
\mcoption{C}{The maximum kinetic energy of
emitted electrons divided by the electron charge}
\mcoption{D}{The intensity of the light}
\ansref{Ch27-A7}

\vspace{0.4cm}

\textbf{A8.} Which of the following provides
evidence that light has particle properties?

\mcoption{A}{Young's double slit experiment}
\mcoption{B}{Diffraction of light through a slit}
\mcoption{C}{The photoelectric effect}
\mcoption{D}{Refraction of light through glass}
\ansref{Ch27-A8}

\vspace{0.4cm}

\textbf{A9.} A hydrogen atom emitting light in
the Lyman series produces photons ending on
the level:

\mcoption{A}{$n=1$}
\mcoption{B}{$n=2$}
\mcoption{C}{$n=3$}
\mcoption{D}{$n=4$}
\ansref{Ch27-A9}

\vspace{0.4cm}

\textbf{A10.} The Heisenberg Uncertainty Principle
states that:

\mcoption{A}{We cannot measure position and
momentum simultaneously with arbitrary precision,
due to instrument limitations}
\mcoption{B}{A particle does not simultaneously
possess definite position and definite momentum ---
this is a fundamental property of nature}
\mcoption{C}{Speed and frequency cannot be
measured simultaneously}
\mcoption{D}{Only applies to photons, not to
massive particles}
\ansref{Ch27-A10}

\vspace{0.6cm}

\subsection*{Section B: Theory and Calculations}

\textbf{B1.} In Rutherford's gold-foil experiment,
a beam of alpha particles ($q = +2e$,
$m = 6.64\times10^{-27}\ \text{kg}$) with kinetic
energy $7.7\ \text{MeV}$ is fired at gold nuclei
($Z = 79$).\\[4pt]
(a) Convert $7.7\ \text{MeV}$ to joules.\\
(b) The distance of closest approach $d$ for a
head-on collision is found from
$KE = kZq_\alpha q_{Au}/d$.
($k = 9.0\times10^9\ \text{N\,m}^2\text{C}^{-2}$,
$q_{Au} = Ze = 79 \times 1.6\times10^{-19}\ \text{C}$)
Calculate $d$ and compare with the nuclear radius
($\sim 7\times10^{-15}\ \text{m}$).\\
(c) Why do most alpha particles pass straight
through the gold foil? Give a quantitative
argument using atomic and nuclear radii.\\
(d) Explain what is meant by the statement that
the atom is ``mostly empty space.''
\hfill\ansref{Ch27-B1}

\vspace{0.4cm}

\textbf{B2.} A hydrogen atom in the ground state
absorbs a photon and is excited.\\
($h = 6.63\times10^{-34}\ \text{J\,s}$,
$c = 3.0\times10^8\ \text{m\,s}^{-1}$,
$1\ \text{eV} = 1.6\times10^{-19}\ \text{J}$)\\[4pt]
(a) Calculate the energy in eV of a photon of
wavelength $97.2\ \text{nm}$ (ultraviolet).\\
(b) Which Bohr transition does this photon cause?
(Use $E_n = -13.6/n^2$ eV)\\
(c) Calculate the wavelength of the photon emitted
when the electron returns from $n=4$ to $n=2$.\\
(d) Is this wavelength in the visible range? What
colour would it appear?\\
(e) Calculate the wavelength of the photon that
would just ionise a hydrogen atom from the ground
state.
\hfill\ansref{Ch27-B2}

\vspace{0.4cm}

\textbf{B3.} A zinc photoelectric cell has work
function $\phi = 4.3\ \text{eV}$.
($h = 6.63\times10^{-34}\ \text{J\,s}$,
$c = 3.0\times10^8\ \text{m\,s}^{-1}$,
$m_e = 9.11\times10^{-31}\ \text{kg}$,
$e = 1.6\times10^{-19}\ \text{C}$)\\[4pt]
(a) Calculate the threshold frequency and
threshold wavelength for zinc.\\
(b) A lamp emits UV light of wavelength $200\ \text{nm}$.
Calculate the maximum kinetic energy and maximum
speed of photoelectrons.\\
(c) Calculate the stopping potential.\\
(d) The lamp's intensity is doubled. State and
explain what happens to: (i) the maximum kinetic
energy; (ii) the rate of electron emission.\\
(e) A second metal has threshold wavelength
$500\ \text{nm}$ (green light). Calculate its
work function. Which requires higher-energy photons
to cause emission, zinc or this metal?
\hfill\ansref{Ch27-B3}

\vspace{0.4cm}

\textbf{B4.} An electron is accelerated from rest
through a potential difference of $V$.\\
($h = 6.63\times10^{-34}\ \text{J\,s}$,
$m_e = 9.11\times10^{-31}\ \text{kg}$,
$e = 1.6\times10^{-19}\ \text{C}$)\\[4pt]
(a) Show that the de Broglie wavelength of the
accelerated electron is:
\[
  \lambda = \frac{h}{\sqrt{2m_e eV}}
\]
(b) Calculate $\lambda$ for $V = 50\ \text{V}$,
$200\ \text{V}$, and $10\,000\ \text{V}$.\\
(c) Typical crystal lattice spacings are
$0.1$--$0.5\ \text{nm}$. For which of these voltages
would you expect electron diffraction?\\
(d) Compare your wavelengths in (b) with the
wavelength of visible light ($400$--$700\ \text{nm}$)
and explain why electron microscopes can resolve
features far smaller than optical microscopes.\\
(e) A proton is accelerated through $200\ \text{V}$.
Calculate its de Broglie wavelength and compare
with the electron at the same voltage.
($m_p = 1.67\times10^{-27}\ \text{kg}$)
\hfill\ansref{Ch27-B4}

\begin{center}
\begin{tikzpicture}[font=\small]
  \begin{axis}[
    xmode=log,log basis x=10,
    xmin=0.001,xmax=1000,ymin=0,ymax=1.45,
    width=12cm,height=3.2cm,
    xtick={0.001,0.01,0.1,1,10,100,1000},
    xticklabels={$10^{-3}$,$10^{-2}$,$10^{-1}$,$1$,$10$,$10^2$,$10^3$},
    ytick=\empty,axis lines=left,
    xlabel={wavelength $\lambda$ (nm)},
    tick label style={font=\footnotesize},
    label style={font=\small},
    clip=false
  ]
    \addplot[draw=none,fill=codexblue!22] coordinates
      {(0.1,0.38) (0.5,0.38) (0.5,0.82) (0.1,0.82)} -- cycle;
    \addplot[draw=none,fill=codexgold!30] coordinates
      {(400,0.38) (700,0.38) (700,0.82) (400,0.82)} -- cycle;
    \node[align=center,font=\footnotesize,text width=2.8cm] at (axis cs:0.22,1.05)
      {crystal spacing: $0.1$--$0.5\,\text{nm}$};
    \node[align=center,font=\footnotesize,text width=2.8cm] at (axis cs:530,1.05)
      {visible light: $400$--$700\,\text{nm}$};
  \end{axis}
\end{tikzpicture}
\par\smallskip\footnotesize Place your calculated wavelengths on this logarithmic scale.
\end{center}

\vspace{0.4cm}

\textbf{B5.} This question compares classical and
quantum descriptions of the hydrogen atom.\\[4pt]
(a) In the Bohr model, the electron in the ground
state orbits at radius $r_1 = 5.29\times10^{-11}\ \text{m}$.
Using the centripetal force and Coulomb force:
\[
  \frac{m_e v^2}{r} = \frac{ke^2}{r^2}
\]
calculate the speed of the electron in the ground state.\\
(b) Calculate the de Broglie wavelength of this
electron and compare with the circumference of
the Bohr orbit ($2\pi r_1$). What do you notice?\\
(c) The classical model predicts the electron
should radiate and spiral into the nucleus in
about $10^{-11}\ \text{s}$. Estimate the rate of
energy loss (power radiated) using the classical
formula for an accelerating charge:
$P = e^2a^2/(6\pi\varepsilon_0 c^3)$,
given the electron acceleration is
$9.0\times10^{22}\ \text{m\,s}^{-2}$.
Comment on why atoms are actually stable.\\
(d) Explain in your own words why the Bohr model,
despite correctly predicting hydrogen spectral
lines, was known to be incomplete even when
Bohr proposed it.
\hfill\ansref{Ch27-B5}

\begin{center}
\begin{tikzpicture}[scale=0.9,transform shape,font=\small]
  % Classical prediction: radiating orbit contracts toward the nucleus.
  \node[font=\bfseries] at (1.8,2.7) {Classical model};
  \fill[codexred] (1.9,1.35) circle (0.12);
  \draw[dashed,gray] (1.9,1.35) circle (0.95);
  \draw[->,thick,codexblue] (2.85,1.35) arc[start angle=0,end angle=300,radius=0.95];
  \fill[codexblue] (2.38,0.55) circle (0.08);
  \draw[->,thick,codexred] (2.38,0.55) -- (2.02,1.05);
  % A diffuse stationary-state density, not a fixed classical orbit.
  \node[font=\bfseries] at (6.4,2.7) {Quantum state};
  \shade[inner color=codexblue!35,outer color=codexblue!4]
    (6.4,1.35) circle (1.05);
  \fill[codexred] (6.4,1.35) circle (0.12);
  \draw[gray,dashed] (4.0,0.0) -- (4.0,2.9);
\end{tikzpicture}
\par\smallskip\footnotesize The diffuse cloud represents probability, not a fixed orbit.
\end{center}

\end{exercisebox}
